\section{Data Feed - Anpassung}
%
\subsection{Primäre Personas}
%
Benjamin, Leah und Ralph nutzen die App sehr aktiv für ihr Hobby, Bildung und Beruf. Einen lockereren und einfachen Ansatz zur Datendarstellung werden sie sicher wertschätzen, aber 
sie wollen vermutlich mehr. Um "Power Usern" ein besseres Erlebnis zu bieten, wird eine Suchfunktion mit Filtern eingebaut. Standardmäßig kann man mit dieser das aktuell gewählte
Thema durchsuchen, es gibt aber auch weitere Optionen wie Filter, Quellen und ob man z.B.: Den Text oder nur den Titel durchsuchen will.
%
\begin{figure}[H]
    \centering
    \includegraphics[width=0.90\textwidth]{app_erweiterung_primaer_alexander.png}
    \caption{Erweiterung für Primäre Nutzer}
    \label{fig:Primaere_Erweiterung_Alexander}
\end{figure}
%
\noindent
In die rechte obere Ecke der Hauptseite wird ein Lupensymbol hinzugefügt. Wird dieses angeklickt, dann wird die Suchleiste eingeblendet und der Hintergrund unscharf gemacht.
Die Tastatur öffnet sich automatisch und man kann sofort tippen. Durch Klicken außerhalb der Suchleiste wird die Suche beendet.
Wird das Burgermenü angeklickt, dann öffnen sich die weiteren Einstellungen. Hier können die Suchparameter angepasst werden.
\\
\\
Bei jedem Post wird an die rechte obere Ecke ein "Lesezeichen"-Symbol platziert. Hiermit können die Nutzer den Post für später speichern.
%
\subsection{Sekundäre Personas}
%
Luca, unser Repräsentant der sekundären Personas, ist ein Nutzer, der die App nicht oft nutzt. Während der Nutzung gibt es zwar ein Ziel, aber innerhalb der App bei unseren Daten
wird nichts spezifisch genutzt. Für die sekundären Nutzer werden zwei Dinge hinzugefügt. Erstens wird noch vor der Auswahl der Themen ein weiterer Unterpunkt hinzugefügt: die Daten des Tages/Woche. 
Hier wird den Nutzern ein aktueller und interessanter Post vorgeschlagen. Zusätzlich gibt es noch eine Zufallssuche, hier kann man auf Knopfdruck einen zufälligen Post heraussuchen.
%
\begin{figure}[H]
    \centering
    \includegraphics[width=0.30\textwidth]{app_erweiterung_sekundaer_alexander.png}
    \caption{Erweiterung für Sekundäre Nutzer}
    \label{fig:Sekundaere_Erweiterung_Alexander}
\end{figure}
%
\subsection{Negative Persona}
%
Die negative Persona ist jemand, der die dargestellten Daten zweckentfremden will und damit die Meinung anderer manipulieren will. In einer weniger extremen Darstellung
eines negativen Nutzers könnte doch auch jemand versuchen, die Daten entweder falsch auszulegen oder sogar die Daten selber zu verändern.
Innerhalb der App / der UI gibt es hier weniger zu tun. Eher ist es wichtiger, die API und App selbst mit guten Sicherheitsgrundlagen zu programmieren. 
Zusätzlich müssen die rechtlichen Grundlagen gut geklärt sein, wo und wie die Daten weiterverwendet werden dürfen (Lizenzen etc.).
Seitens der Quellen muss es auch klar sein, dass wir keine Verantwortung für die Richtigkeit oder auch den Inhalt der Datensätze übernehmen.